\begin{center}
    {\LARGE \textbf{2015无机化学}} \\[2ex]
    {\large 时量180分钟 \quad 满分100分}
\end{center}

\vspace{0.5cm}
\noindent\textbf{注意事项:}
\begin{enumerate}[label=\arabic*., parsep=0pt, itemsep=0pt]
    \item 答题前,考生先将自己的姓名、学号、考场号、座位号填写在答题卡上,将条形码准确粘贴在考生信息条形码粘贴区。
    \item 考试结束后,将本试卷与答题纸一并收回。
    \item 回答各个问题时,将答案写在答题纸上,写在本试卷上无效。
\end{enumerate}

\vspace{0.5cm}
\noindent\textbf{一、完成并配平下列化学反应方程式：（15 pts.）}

\begin{enumerate}[label=\arabic*., start=1, itemsep=1.5ex]
    \item \ce{(NH4)2Cr2O7}受热分解。
    \item 白磷与氢氧化钠溶液作用。
    \item 向\ce{FeCl3}溶液中通入\ce{H2S}。
    \item 将次氯酸钠的碱性溶液滴入到硝酸铅溶液中。
    \item 将\ce{SeO2}溶于水，然后通入\ce{SO2}气体。
\end{enumerate}

\vspace{0.5cm}
\noindent\textbf{二、解释下列现象，写出必要的反应方程式。（25 pts.）}

\begin{enumerate}[label=\arabic*., start=1, itemsep=1.5ex]
    \item 打开盛装\ce{TiCl4}试剂的玻璃瓶时会冒白烟。
    \item \ce{Fe^{3+}}离子可以被\ce{I^-}离子还原为\ce{Fe^{2+}}离子，并生成单质\ce{I2}，但在\ce{Fe^{3+}}离子溶液中先加入一定量氟化物，再加入\ce{I^-}离子，则不会有\ce{I2}生成。
    \item 将\ce{SO2}通入到\ce{CuSO4}和\ce{NaCl}的浓的混合溶液中，有白色沉淀生成。
    \item 向酸性\ce{Cr2O7^{2-}}溶液中加入金属锌，溶液的颜色由橙色变为绿色，继而变成蓝色，放置后又变回绿色。
    \item 将少量酸性\ce{MnSO4}溶液与\ce{(NH4)2S2O8}溶液混合后水浴加热，很快生成棕黑色沉淀；但在加热前加入几滴\ce{AgNO3}溶液，混合溶液逐渐变红。
\end{enumerate}

\vspace{0.5cm}
\noindent\textbf{三、回答下列问题（70 pts.）}

\begin{enumerate}[label=\arabic*., start=1, itemsep=1.5ex]
    \item 写出下列所有元素的元素符号和名称：
    \begin{enumerate}[label=(\arabic*), itemsep=1ex]
        \item 具有价电子构型为$(n-1)d^{10}ns^1$的所有元素；
        \item $4p$轨道上含有单电子的所有元素。
    \end{enumerate}
    \item 比较大小，并简要说明理由。
    \begin{enumerate}[label=(\arabic*), itemsep=1ex]
        \item 键角：$\ce{NO2^+}$，$\ce{NO2^-}$，$\ce{NO3^-}$
        \item 稳定性：$\ce{OF}$，$\ce{OF^+}$，$\ce{OF^-}$
    \end{enumerate}
    \item 试比较\ce{NH3}，三甲胺$\ce{N(CH3)3}$，三甲硅胺$\ce{N(SiH3)3}$的分子结构哪一个明显不同，为什么？
    \item 已知$\ce{I3^-}$和$\ce{N3^-}$均具有直线形结构，试用杂化轨道理论说明两者的杂化方式及成键情况。
    \item 比较氯的四种含氧酸：$\ce{HClO}$，$\ce{HClO2}$，$\ce{HClO3}$和$\ce{HClO4}$的酸性强弱，并解释原因。
    \item 将\ce{NH3}和$\ce{AsH3}$分别通入\ce{AgNO3}溶液中时，现象及产物是否相同？为什么？写出反应方程式。
    \item 试说明下列问题中加（浓）盐酸的作用是什么？写出相关的反应方程式。
    \begin{enumerate}[label=(\arabic*), itemsep=1ex]
        \item \ce{Na3AsO4}与\ce{KI}溶液反应。
        \item 单质金与浓硝酸反应。
        \item 配置\ce{SnCl2}溶液。
    \end{enumerate}
    \item 为什么$\ce{K4[Fe(CN)6] * 3H2O}$可由$\ce{FeSO4}$溶液与$\ce{KCN}$直接混合来制备，而$\ce{K3[Fe(CN)6]}$却不能由$\ce{FeCl3}$溶液与$\ce{KCN}$直接混合来制备？如何制备赤血盐？
    \item 比较下列配位化合物的稳定性，简要说明理由。
    \begin{enumerate}[label=(\arabic*), itemsep=1ex]
        \item $\ce{[HgI4]^{2-}}$ 与 $\ce{[HgCl4]^{2-}}$
        \item $\ce{[Cu(en)2]^{2+}}$ 与 $\ce{[Cu(NH3)4]^{2+}}$
        \item $\ce{[Hg(CN)4]^{2-}}$ 与 $\ce{[Cd(CN)4]^{2-}}$
        \item $\ce{[Ag(NH3)2]^+}$ 与 $\ce{[Ag(S2O3)2]^{3-}}$
    \end{enumerate}
    \item 测得电池：
    
    $\ce{Pb(s)|Pb^{2+}(10^{-3}\,mol\cdot dm^{-3})\parallel VO^{2+}(10^{-3}\,mol\cdot dm^{-3})|}$\\
    $\ce{V^{3+}(10^{-5}\,mol\cdot dm^{-3})}$，$\ce{H^+}(10^{-3}\,mol\cdot dm^{-3})|\ce{Pt(s)}$
    
    的电动势为$+0.67V$。计算：
    \begin{enumerate}[label=(\arabic*), itemsep=1ex]
        \item 电对$\ce{VO^{2+}/V^{3+}}$的$E^\ominus$；
        \item 298K时，反应$\ce{Pb(s) + 2VO^{2+} + 4H^+ = Pb^{2+} + 2V^{3+} + 2H2O}$的平衡常数$K^\ominus$。
    \end{enumerate}
    （已知：$E^\ominus(\ce{Pb^{2+}/Pb})=-0.126V$）
\end{enumerate}

\vspace{0.5cm}
\noindent\textbf{四、分离、鉴别与制备（20 pts.）}

\begin{enumerate}[label=\arabic*., start=1, itemsep=1.5ex]
    \item 以$\ce{Na2S}$、$\ce{Na2CO3}$和$\ce{SO2}$为原料制备硫代硫酸钠。
    \item \ce{CuCl}，\ce{AgCl}，\ce{Hg2Cl2}均为难溶于水的白色粉末，试区别这三种物质。
    \item 有一黄色固体，可能是$\ce{PbCrO4}$或$\ce{BaCrO4}$，试加以鉴别。
    \item 设计方案分离混合溶液中的离子：$\ce{Fe^{3+}}$、$\ce{Al^{3+}}$、$\ce{Zn^{2+}}$、$\ce{Ag^+}$
\end{enumerate}

\vspace{0.5cm}
\noindent\textbf{五、推断（20 pts.）}

\begin{enumerate}[label=\arabic*., start=1, itemsep=1.5ex]
    \item 某蓝色固体$\ce{A}$易分解为无色气体$\ce{B}$和红棕色气体$\ce{C}$。气体$\ce{B}$高温显顺磁性，低温显逆磁性。它在常温下易在空气中变为红棕色气体$\ce{C}$。气体$\ce{B}$、$\ce{C}$均有很高的毒性。除去$\ce{B}$的简便方法是用硫酸亚铁溶液吸收，这是因为能生成配合物$\ce{D}$的缘故。$\ce{C}$被$\ce{NH3}$在一定温度下经催化还原为无毒的气体$\ce{E}$。$\ce{A}$溶于水变成天蓝色无机酸$\ce{F}$，$\ce{F}$易歧化为$\ce{G}$酸和气体$\ce{B}$；$\ce{B}$和氯气以$2:1$（摩尔比）直接作用，生成橙色气体$\ce{H}$；无机酸$\ce{F}$与尿素快速反应，生成气体$\ce{E}$和气体$\ce{I}$；气体$\ce{C}$易聚合得无色化合物$\ce{J}$。金属铜可溶解在$\ce{J}$中。根据以上实验现象，指出$\ce{A\rightarrow J}$所代表的物质的化学式，给出相关各步的化学反应方程式。
    \item 有一固体混合物可能含有$\ce{AlCl3}$，$\ce{KMnO4}$，$\ce{K2SO4}$，$\ce{AgNO3}$，$\ce{CuS}$和$\ce{ZnCl2}$中的一种或几种。将此混合物置于水中，并用少量盐酸酸化，得白色沉淀物$\ce{A}$和无色溶液$\ce{B}$。白色沉淀$\ce{A}$可溶于氨水中。将滤液$\ce{B}$分成两份，一份中加入少量$\ce{NaOH}$溶液，有白色沉淀产生，再加入过量$\ce{NaOH}$溶液则白色沉淀溶解。另一份中加入少量氨水，也产生白色沉淀，当加入过量氨水时，白色沉淀溶解。
    试根据上述现象，判断在混合物中，哪些物质肯定存在？哪些肯定不存在？哪些可能存在？简述理由，并用方程式表示。
\end{enumerate}